% TOP_PROBLEM: {"id": 2305039, "title": "Research Problems in Function Theory — Problem 5.39", "category_id": 8, "category": {"id": 8, "name": "analysis", "display_name": "Analysis", "description": "Limits, continuity, calculus, and function theory.", "slug": "analysis", "order_index": 8, "created_at": "2026-07-31T15:26:25.671Z"}, "status": "open", "problem_number": "AMR-022-5039", "set_id": 13}
% FIRST_POSTED: 2026-10-01
% ENTRY_KIND: statement_only
% UnsolvedMath ID 2305039; source catalogue code AMR-022-5039
% !TeX program = lualatex
% Definition and sources only; this file is not an LLM solution attempt.
\documentclass[11pt,a4paper]{article}
\usepackage[margin=25mm]{geometry}
\usepackage{fontspec}
\setmainfont{lmroman10-regular.otf}[BoldFont=lmroman10-bold.otf,
 ItalicFont=lmroman10-italic.otf,BoldItalicFont=lmroman10-bolditalic.otf]
\usepackage{amsmath,amssymb,mathtools}
\usepackage{unicode-math}
\setmathfont{latinmodern-math.otf}
\AtBeginDocument{\let\setminus\smallsetminus}
\usepackage{xurl,hyperref}
\hypersetup{colorlinks=true,urlcolor=blue}
\setcounter{secnumdepth}{0}
\setlength{\parindent}{0pt}
\setlength{\parskip}{5pt}
\setlength{\emergencystretch}{3em}
\begin{document}
\section*{Research Problems in Function Theory — Problem 5.39}\label{2305039}
{\small UnsolvedMath ID 2305039 · AMR-022-5039 · Analysis · AMR Open Problem Lists}
\subsection{Definitions and mathematical statement}
For $p>0$ and an analytic function $h$ on the unit disk
$\mathbb D$, define its radial $p$-mean by
\[
 M_p(r,h)=\left(\frac1{2\pi}\int_0^{2\pi}
                     |h(re^{i\theta})|^p\,d\theta\right)^{1/p}
 \quad(0\leq r<1).
\]
An analytic function $g$ is subordinate to $f$ if
$g=f\circ\phi$ for an analytic map $\phi:\mathbb D\to\mathbb D$
with $\phi(0)=0$. No injectivity assumption is imposed on $f$.

Let $r_p$ be the supremum of all $\rho\in[0,1]$ for which
\[
 M_p(r,(f\circ\phi)')\leq M_p(r,f')
 \quad\text{for every }0<r<\rho
\]
holds for every analytic $f$ and every such map $\phi$.
Determine $r_p$ exactly as a function of $p$, and establish
sharpness by examples that violate the inequality beyond the
optimal radius. The universal quantifiers include constant
functions, which cause no difficulty since their derivatives vanish.

The comparison concerns derivatives; subordination of the functions
themselves does not automatically give this derivative inequality
throughout the disk. The simple choice $f(z)=z$ and $\phi(z)=z^2$
makes the two derivative means $1$ and $2r$, respectively,
showing that $r_p\leq1/2$ for every $p$. The classical case
$p=2$ has radius $1/2$; the question asks for the full range
of positive exponents, including $0<p<1$.
\subsection{Short English statement}
Find the largest universal radius on which subordination also decreases the radial mean of the derivative.
\subsection{Sources}
\begin{itemize}
\item[S1] ULAM.AI and the UnsolvedMath Contributors, supplied problems.json, record 2305039 (AMR-022-5039), AMR Open Problem Lists. Catalogue identity and source transcription; CC BY 4.0.
\url{https://huggingface.co/datasets/ulamai/UnsolvedMath}.
\item[S2] Hayman - Research Problems in Function Theory (2018), Problem 5.39. Original source indexed by AMR-022-5039.
\url{https://arxiv.org/abs/1809.07200}.
\end{itemize}
\end{document}
